\section{Conditioning on a tested set}\label{sec:exposure}

Fix a finite-dimensional strategy for the $k$-fold game, with state
$\psi$ and POVMs $\{M_{\mathbf x}^{\mathbf a}\}$ and
$\{N_{\mathbf y}^{\mathbf b}\}$. Let $C\subseteq[k]$ be a set of tested
coordinates and let
\[
 E_C=\{V(X_j,Y_j,A_j,B_j)=1\text{ for every }j\in C\},
 \qquad p=\Pr(E_C)>0.
\]
Assume that $m=k-|C|\ge1$, and put
\begin{equation}\label{eq:conditioning-budget}
 d_C=(|\mathcal A||\mathcal B|)^{|C|},\qquad
 L=\log(d_C/p),\qquad \tau=L/m.
\end{equation}
We construct a vector and local measurements that reproduce the
conditional answers at an untested coordinate. Two approximating vectors,
one determined by each player's question and a transcript, will have
total squared error at most $2p\tau$ on average. These comparisons let
Section~\ref{app:root-rounding} replace the conditional vector by a fixed
state on each support component. Throughout this
section, conditional expectations of operators use the original,
unconditioned question distribution $\mu^{\otimes k}$. We adapt the
weighted reveal argument of \cite[Chapter~6, Section~4]{tenadvances2026},
using the fixed ordering below.

\subsection{Exposing one question at each coordinate}

Fix any ordering of the untested coordinates, and abbreviate their
questions as $X_1,\ldots,X_m$ and $Y_1,\ldots,Y_m$, with answers
$A_1,\ldots,A_m$ and $B_1,\ldots,B_m$. The tested questions
retain the notation $X_C,Y_C$. Position-indexed win indicators refer to
these same untested coordinates; sums over $i\notin C$ in the result
statements use the original coordinate labels. For $0\le j\le m$
define the flag
\[
 \Omega_j=(X_C,Y_C,X_{\le j},Y_{>j}).
\]
Conditioned on a flag of positive probability, the full question vectors
$\mathbf X$ and $\mathbf Y$ are independent. Indeed, at a tested
coordinate both questions are fixed. At each untested coordinate one
question is fixed and the other has its conditional law under $\mu$;
these remaining choices are independent across coordinates. This
argument uses no independence between the two questions within a
coordinate.

For tested answer strings $\alpha\in\mathcal A^C$ and
$\gamma\in\mathcal B^C$, define the coarse-grained effects
\[
 M_{\mathbf x}^{\alpha}
   =\sum_{\mathbf a:\,a_C=\alpha}M_{\mathbf x}^{\mathbf a},
 \qquad
 N_{\mathbf y}^{\gamma}
   =\sum_{\mathbf b:\,b_C=\gamma}N_{\mathbf y}^{\mathbf b},
\]
and the testing mask
\[
 \theta(x_C,y_C,\alpha,\gamma)
   =\prod_{j\in C}V(x_j,y_j,\alpha_j,\gamma_j).
\]
For $C=\varnothing$ there is one pair of empty answer strings and the
empty product is $1$. Write
\[
 F_j^\alpha=\mathbb E[M_{\mathbf X}^\alpha\mid\Omega_j],\qquad
 H_j^\gamma=\mathbb E[N_{\mathbf Y}^\gamma\mid\Omega_j],\qquad
 \chi(A,B)=\langle\psi|A\otimes B|\psi\rangle.
\]
The answer strings here are fixed indices of the effects, not additional
conditions on the question distribution. These are positive contractions.
Conditional independence and linearity
give, for every $j$,
\begin{equation}\label{eq:exposure-mass}
 \mathbb E_{\Omega_j}\sum_{\alpha,\gamma}
       \theta\,\chi(F_j^\alpha,H_j^\gamma)=p.
\end{equation}

A mode is an index $i\in[m]$. Delete its question from the flag and set
\[
 r_0=(X_C,Y_C,X_{<i},Y_{>i}),\qquad \nu_i=\operatorname{law}(r_0).
\]
Figure~\ref{fig:exposure} records which questions are retained.
\Needspace{14\baselineskip}
\begin{figure}[!h]
\centering
\begin{tabular}{c@{\qquad}ccc}
\toprule
 & Before $i$ & Coordinate $i$ & After $i$\\
\midrule
Alice & $X_{<i}$ revealed & $X_i$ fresh & $X_{>i}$ unrevealed\\
Bob   & $Y_{<i}$ unrevealed & $Y_i$ fresh & $Y_{>i}$ revealed\\
\bottomrule
\end{tabular}
\caption{The questions exposed in $r_0$, in addition to $X_C,Y_C$;
$i$ indexes the fixed ordering of the untested coordinates. Under the
original product question law, before conditioning on $E_C$, fixing
$r_0$ and $(X_i,Y_i)$ makes the remaining Alice and Bob question strings
independent. Replacing $Y_i$ by $X_i$ in successive flags produces the
two entropy telescopes.}
\label{fig:exposure}
\end{figure}

The variable $r_0$ is independent of $(X_i,Y_i)\sim\mu$. For each value
of $r_0$ of positive probability, put
\begin{align*}
 F_x^\alpha&=\mathbb E[M_{\mathbf X}^\alpha\mid r_0,X_i=x],&
 H_y^\gamma&=\mathbb E[N_{\mathbf Y}^\gamma\mid r_0,Y_i=y],\\
 F_y^\alpha&=\sum_x\mu(x\mid y)F_x^\alpha,&
 H_x^\gamma&=\sum_y\mu(y\mid x)H_y^\gamma.
\end{align*}
We use these formulas only on positive marginal probabilities. In each
conditional ensemble, zero weights are omitted before applying
Lemma~\ref{lem:resolvent}. Questions of zero marginal probability may be
assigned arbitrary complete local POVMs with answers in the original
alphabets. All four effects also depend on $i,r_0$, which are suppressed
in the notation.
The subscripts in $F_y^\alpha,H_x^\gamma$ indicate the question on which
the conditional average depends. In particular,
\[
 (F_{i-1}^\alpha,H_{i-1}^\gamma)=(F_y^\alpha,H_y^\gamma),
 \qquad
 (F_i^\alpha,H_i^\gamma)=(F_x^\alpha,H_x^\gamma).
\]
Here the first identity is evaluated at $Y_i=y$, and the second at
$X_i=x$. The mask depends only on $r_0,\alpha,\gamma$, not on $x,y$.

\subsection{The conditional vector and its measurements}

Apply the common finite-dimensional compression in
Lemma~\ref{lem:resolvent} to all four families
$F_x^\alpha,F_y^\alpha,H_x^\gamma,H_y^\gamma$, over every mode,
flag and positive-marginal question. This gives fixed finite-dimensional
target spaces $\mathcal K_A,\mathcal K_B$ for the maps $\mathcal F_T$,
including the maps used later for the approximations.

Let $R_i$ be the finite set of transcripts $r=(r_0,\alpha,\gamma)$,
where $r_0$ ranges over its positive-probability values and
$(\alpha,\gamma)\in\mathcal A^C\times\mathcal B^C$. Put
\begin{gather*}
 w_i(r)=\nu_i(r_0)\theta(r),\qquad
 \Psi_{xy}^r=(\mathcal F_{F_x^\alpha}\otimes
                         \mathcal F_{H_y^\gamma})\psi,\\
 p_{xy}^r=\|\Psi_{xy}^r\|^2=\chi(F_x^\alpha,H_y^\gamma).
\end{gather*}
The last identity follows from $\mathcal F_T^*\mathcal F_T=T$.
The weights $w_i(r)$ use counting measure on the answer strings; they
are independent of $(x,y)$ and need not sum to one. Their sum is at
most $d_C$. The vector's squared norm supplies the remaining probability
of the tested answers.

We next exhibit measurements on $\Psi_{xy}^r$ reproducing the
conditional answers. Define the finer effects
\[
 G_x^{\alpha,a}
  =\mathbb E[M_{\mathbf X}^{\alpha,a_i=a}\mid r_0,X_i=x],\qquad
 K_y^{\gamma,b}
  =\mathbb E[N_{\mathbf Y}^{\gamma,b_i=b}\mid r_0,Y_i=y],
\]
where an unmentioned answer is summed out. They are positive and satisfy
$\sum_aG_x^{\alpha,a}=F_x^\alpha$ and
$\sum_bK_y^{\gamma,b}=H_y^\gamma$. Lemma~\ref{lem:measurement-transport}
provides POVMs $\{\widetilde M_{r,x}^a\}$ on $\mathcal K_A$ and
$\{\widetilde N_{r,y}^b\}$ on $\mathcal K_B$ such that
\[
 \mathcal F_{F_x^\alpha}^*\widetilde M_{r,x}^a
      \mathcal F_{F_x^\alpha}=G_x^{\alpha,a},\qquad
 \mathcal F_{H_y^\gamma}^*\widetilde N_{r,y}^b
      \mathcal F_{H_y^\gamma}=K_y^{\gamma,b}.
\]
If $p_{xy}^r>0$, write
$\varphi_{xy}^r=\Psi_{xy}^r/\sqrt{p_{xy}^r}$. Then
\begin{equation}\label{eq:exposure-fine-law}
 \langle\varphi_{xy}^r|
      \widetilde M_{r,x}^a\otimes\widetilde N_{r,y}^b
       |\varphi_{xy}^r\rangle
  =\frac{\chi(G_x^{\alpha,a},K_y^{\gamma,b})}{p_{xy}^r}.
\end{equation}
To identify this probability, condition the original experiment on
$r_0,X_i=x,Y_i=y$. The unfixed Alice questions $X_{>i}$ and Bob
questions $Y_{<i}$ are independent, with conditional laws determined
by $Y_{>i}$ and $X_{<i}$ respectively. More explicitly, for positive-mass
conditioning data their joint law is
\[
 \Pr(X_{>i}=\mathbf u,Y_{<i}=\mathbf z\mid r_0,X_i=x,Y_i=y)
 =\prod_{j>i}\mu(u_j\mid y_j)
  \prod_{j<i}\mu(z_j\mid x_j).
\]
This is a statement about the original question law, before conditioning
on tested answers. Neither product depends on the fresh pair $(x,y)$.
Since Alice's effect depends only on $\mathbf X$, dropping the condition
$Y_i=y$ leaves its average unchanged; the analogous assertion holds for
Bob after dropping $X_i=x$. Hence averaging the joint effect
$M_{\mathbf X}^{\alpha,a_i=a}\otimes
 N_{\mathbf Y}^{\gamma,b_i=b}$ gives exactly
$G_x^{\alpha,a}\otimes K_y^{\gamma,b}$. Its coarse-grained probability
is $p_{xy}^r$. When $\theta(r)=1$, conditioning further on the tested
answers $(\alpha,\gamma)$ already implies $E_C$. Thus
\eqref{eq:exposure-fine-law} is the conditional distribution of the
untested answers $(A_i,B_i)$ given $E_C,r,X_i=x,Y_i=y$. Values with
$\theta(r)p_{xy}^r=0$ have no mass in that conditional experiment.

In particular, the joint law of the transcript and fresh questions
conditioned on $E_C$ is
\begin{equation}\label{eq:exposure-law}
 \lambda_i(r,x,y)=\frac{w_i(r)\mu(x,y)\|\Psi_{xy}^r\|^2}{p}.
\end{equation}
Thus $\lambda_i$ is a probability law even though the weights $w_i$
alone are not. The measurements in \eqref{eq:exposure-fine-law} are
local: Alice's depends only on $(i,r,x)$ and Bob's only on $(i,r,y)$.
The transcript is an analysis variable, not information given to the
players. A legal one-game strategy must start with a shared resource
independent of the fresh questions and obtain any transcript information
using only that resource and the player's own question.

\subsection{Local descriptions and their approximation error}

The conditional vector $\Psi_{xy}^r$ depends on both fresh questions.
Replace the effect on the opposite player's space by its conditional
average to obtain two vectors with one-question descriptions:
\[
 \Psi_x^r=(\mathcal F_{F_x^\alpha}\otimes
                         \mathcal F_{H_x^\gamma})\psi,\qquad
 \Psi_y^r=(\mathcal F_{F_y^\alpha}\otimes
                         \mathcal F_{H_y^\gamma})\psi.
\]
The first is determined by $(i,r,x)$ and the second by $(i,r,y)$.
Both lie in the same fixed space $\mathcal K_A\otimes\mathcal K_B$
as $\Psi_{xy}^r$. Put $p_x^r=\|\Psi_x^r\|^2$ and
$p_y^r=\|\Psi_y^r\|^2$. The identity
$\mathcal F_T^*\mathcal F_T=T$ and linearity of $\chi$ give
\begin{equation}\label{eq:exposure-norms}
 \begin{gathered}
 p_{xy}^r=\chi(F_x^\alpha,H_y^\gamma),\qquad
 p_x^r=\sum_y\mu(y\mid x)p_{xy}^r,\qquad
 p_y^r=\sum_x\mu(x\mid y)p_{xy}^r,\\
 \mathbb E_{\nu_i\mu}\sum_{\alpha,\gamma}\theta p_{xy}^r
 =\mathbb E_{\nu_i\mu}\sum_{\alpha,\gamma}\theta p_x^r
 =\mathbb E_{\nu_i\mu}\sum_{\alpha,\gamma}\theta p_y^r=p.
 \end{gathered}
\end{equation}
Here the total mass is $p$ by \eqref{eq:exposure-law}; the two averaging
identities give the same mass for the local descriptions.

Define the errors at mode $i$ by
\begin{equation}\label{eq:mode-error-definitions}
 e_A(i)=\mathbb E_{\nu_i\mu}\sum_{\alpha,\gamma}\theta
          \|\Psi_{xy}^r-\Psi_y^r\|^2,\qquad
 e_B(i)=\mathbb E_{\nu_i\mu}\sum_{\alpha,\gamma}\theta
          \|\Psi_{xy}^r-\Psi_x^r\|^2.
\end{equation}
The subscript names the tensor factor that changes: $e_A$ replaces
Alice's effect and compares with $\Psi_y$, which Bob can describe from
$(i,r,y)$; $e_B$ replaces Bob's effect and compares with $\Psi_x$, which
Alice can describe from $(i,r,x)$. We will bound their averages by an
entropy telescope.

\begin{lemma}[Mixed exposure]\label{lem:exposure}
For any finite-dimensional repeated-game strategy, proper tested set
$C$ with $p=\Pr(E_C)>0$, and fixed ordering of the $m=k-|C|$ untested
coordinates, the preceding construction gives finite-dimensional vectors
and local POVMs satisfying the exact conditional answer law
\eqref{eq:exposure-fine-law}, transcript law \eqref{eq:exposure-law},
and norm identities \eqref{eq:exposure-norms}. With
$\tau=(|C|\log(|\mathcal A||\mathcal B|)+\log(1/p))/m$,
\begin{equation}\label{eq:exposure-errors}
 \mathbb E_i e_A(i)\le p\tau,\qquad
 \mathbb E_i e_B(i)\le p\tau,\qquad
 \mathbb E_i\bigl(e_A(i)+e_B(i)\bigr)\le2p\tau,
\end{equation}
where $i$ is uniform on $[m]$ and the errors are defined in
\eqref{eq:mode-error-definitions}. The question distribution may have
zero entries and disconnected support.
\end{lemma}

\begin{proof}
The construction and exact laws have been verified above. To prove the
error bounds, let $f(t)=t\log t$, with $f(0)=0$, as in
Lemma~\ref{lem:resolvent}. Use the flags $\Omega_j$ to define
\[
 U_j=\mathbb E_{\Omega_j}\sum_{\alpha,\gamma}
      \theta\,\chi(f(F_j^\alpha),H_j^\gamma),\qquad
 Z_j=\mathbb E_{\Omega_j}\sum_{\alpha,\gamma}
      \theta\,\chi(F_j^\alpha,f(H_j^\gamma)).
\]
We will show that moving from $\Omega_{i-1}$ to $\Omega_i$ increases
$U$ by at least $e_A(i)$ and decreases $Z$ by at least $e_B(i)$.
The common range bound
\begin{equation}\label{eq:exposure-potential-range}
 -pL\le U_j,Z_j\le0\qquad(0\le j\le m)
\end{equation}
will then bound each sum of errors by $pL$.

For the lower bound on $U_j$, fix a flag and a pair of tested answer
strings. Set $F=F_j^\alpha$, $H=H_j^\gamma$,
$g=\chi(I,H)\le1$, and $h=\chi(F,H)$. The positive functional
$A\mapsto\chi(A,H)$ has mass $g$. Applying scalar concavity of
$h_0(t)=-t\log t$ to a spectral decomposition of $F$ gives
\[
 -\chi(f(F),H)\le g h_0(h/g)
       =h\log(g/h)\le h_0(h).
\]
The formula has value $0$ when $g=0$, since then $h=0$. Because
$\theta\in\{0,1\}$ and $h_0(0)=0$, Jensen's inequality over the flag
and the $d_C$ answer pairs, together with
\eqref{eq:exposure-mass}, gives
\[
 -U_j\le\mathbb E\sum_{\alpha,\gamma}h_0(\theta h)
      \le d_C h_0(p/d_C)=pL.
\]
The proof for $Z_j$ interchanges the two players. The upper bounds in
\eqref{eq:exposure-potential-range} follow from $f(T)\preceq0$ for
positive contractions $T$.

The changes when replacing $Y_i$ by $X_i$ in the flag are
\begin{align}
 J_A(i):=U_i-U_{i-1}
 &=\mathbb E_{\nu_i\mu_Y}\sum_{\alpha,\gamma}\theta\,
   \chi\left(\sum_x\mu(x\mid y)f(F_x^\alpha)
                      -f(F_y^\alpha),H_y^\gamma\right),
       \label{eq:exposure-jensen-A}\\
 J_B(i):=Z_{i-1}-Z_i
 &=\mathbb E_{\nu_i\mu_X}\sum_{\alpha,\gamma}\theta\,
   \chi\left(F_x^\alpha,
       \sum_y\mu(y\mid x)f(H_y^\gamma)-f(H_x^\gamma)\right).
       \label{eq:exposure-jensen-B}
\end{align}
For example, linearity in the second argument rewrites $U_i$ as
$\mathbb E_{\nu_i\mu}\sum_{\alpha,\gamma}
\theta\chi(f(F_x^\alpha),H_y^\gamma)$; subtracting $U_{i-1}$ gives
\eqref{eq:exposure-jensen-A}. The other identity uses linearity in the
first argument.

For fixed $r_0,y,\alpha,\gamma$, apply the Jensen estimate of
Lemma~\ref{lem:resolvent} to the ensemble $F_x^\alpha$ with probabilities
$\mu(x\mid y)$ and then evaluate its operator inequality against
$H_y^\gamma$ in the state $\psi$. This bounds the corresponding
squared distance between $\Psi_{xy}^r$ and $\Psi_y^r$, and hence
$e_A(i)\le J_A(i)$. Applying the same argument to the ensemble
$H_y^\gamma$ with probabilities $\mu(y\mid x)$ gives
$e_B(i)\le J_B(i)$. In particular, both potential changes are
nonnegative. Telescoping along the fixed ordering now yields
\begin{equation}\label{eq:exposure-telescope}
 \frac1m\sum_{i=1}^mJ_A(i)\le p\tau,\qquad
 \frac1m\sum_{i=1}^mJ_B(i)\le p\tau,
\end{equation}
by \eqref{eq:exposure-potential-range} and $L=m\tau$.
These inequalities prove \eqref{eq:exposure-errors}. The ordering was
arbitrary and fixed throughout; no averaging over orderings is required.
\end{proof}

\subsection{The change in the sampled question distribution}

The vectors describe the experiment after conditioning on $E_C$.
A one-game strategy, however, receives questions with the original law
$\mu$. The following estimate bounds the resulting discrepancy. We use
$\|\cdot\|_1$ for the full $\ell_1$ distance of probability distributions.

\begin{lemma}[Question drift]\label{lem:question-drift}
Under the preceding assumptions, put
\[
 \Delta_i=\|\operatorname{law}(X_i,Y_i\mid E_C)-\mu\|_1.
\]
For a uniformly chosen untested coordinate,
\[
 \mathbb E_i\Delta_i
 \le\sqrt{\frac{2\log(1/p)}{m}}\le\sqrt{2\tau}.
\]
\end{lemma}
\begin{proof}
Let $Q=((X_j,Y_j))_{j\in[k]}$ be the full question vector and let
$P=\operatorname{law}(Q\mid E_C)$. On the support of $\mu^{\otimes k}$,
\[
 \frac{P(Q)}{\mu^{\otimes k}(Q)}
   =\frac{\Pr(E_C\mid Q)}p\le\frac1p.
\]
It follows that $D(P\|\mu^{\otimes k})\le\log(1/p)$, where $D$ denotes
relative entropy with natural logarithms. Relative entropy to a product
measure is at least the sum of the relative entropies of its marginals:
this follows by expanding the entropies and using subadditivity.
Retaining only the $m$ untested coordinates, Pinsker's inequality and
Cauchy--Schwarz give
\[
 \frac1m\sum_{i=1}^m\Delta_i
 \le \frac1m\sum_{i=1}^m
       \sqrt{2D(P_i\|\mu)}
 \le\sqrt{\frac2m\sum_{i=1}^mD(P_i\|\mu)}
 \le\sqrt{\frac{2\log(1/p)}m}.
\]
\end{proof}
